% Abstract (~200 words, self-contained, one quantitative result)
\noindent\textbf{Motivation:} Rare cell types in single-cell RNA-seq are systematically missed by semi-supervised
classifiers when only a handful of rare cells are labelled. Simply encouraging more rare-class predictions
recovers some of these cells but inflates false positives, which can be more damaging than the original omissions.

\noindent\textbf{Results:} We present scRareRefine, a post-hoc module that refines the predictions of a frozen
scANVI backbone. It scores each cell by an anisotropic prototype-membership distance computed only on the training
set, and decides which cells to relabel through a conformal threshold and three abstention gates calibrated only on
a validation set, so that the false-rescue rate (FFR) is held below a fixed budget $\alpha=0.01$. Across six datasets,
nine methods and three random splits, scRareRefine raises mean rare-cell F1 in the label-scarce regime from
$0.72$ (scANVI) to $0.88$ while keeping the worst-case FFR at $0.0098\le\alpha$; it is the only method that achieves
both, and it never underperforms scANVI in this regime (paired one-sided Wilcoxon $p=1.3\times10^{-6}$). The method
is conditional by design: it abstains, returning the backbone prediction, whenever rescue is not warranted, and we
characterise its failure modes, including a recall ceiling for transcriptionally entangled rare types.

\noindent\textbf{Availability:} Code and reproduction scripts are available at \url{https://github.com/REPO}.
